\UseRawInputEncoding
\documentclass[a4paper, 14pt]{extarticle}

% Поля
%--------------------------------------
\usepackage{geometry}
\geometry{a4paper,tmargin=2cm,bmargin=2cm,lmargin=3cm,rmargin=1cm}
%--------------------------------------


%Russian-specific packages
%--------------------------------------
\usepackage[T2A]{fontenc}
\usepackage[utf8]{inputenc} 
\usepackage[english, main=russian]{babel}
%--------------------------------------

\usepackage{textcomp}

% Красная строка
%--------------------------------------
\usepackage{indentfirst}               
%--------------------------------------             


%Graphics
%--------------------------------------
\usepackage{graphicx}
\graphicspath{ {./images/} }
\usepackage{wrapfig}
%--------------------------------------

% Полуторный интервал
%--------------------------------------
\linespread{1.3}                    
%--------------------------------------

%Выравнивание и переносы
%--------------------------------------
% Избавляемся от переполнений
\sloppy
% Запрещаем разрыв страницы после первой строки абзаца
\clubpenalty=10000
% Запрещаем разрыв страницы после последней строки абзаца
\widowpenalty=10000
%--------------------------------------

%Списки
\usepackage{enumitem}

%Подписи
\usepackage{caption} 

%Гиперссылки
\usepackage{hyperref}

\hypersetup {
	unicode=true
}

%Рисунки
%--------------------------------------
\DeclareCaptionLabelSeparator*{emdash}{~--- }
\captionsetup[figure]{labelsep=emdash,font=onehalfspacing,position=bottom}
%--------------------------------------

\usepackage{tempora}
\usepackage{amsmath}
\usepackage{color}
\usepackage{listings}
\lstset{
  belowcaptionskip=1\baselineskip,
  breaklines=true,
  frame=L,
  xleftmargin=\parindent,
  language=Python,
  showstringspaces=false,
  basicstyle=\footnotesize\ttfamily,
  keywordstyle=\bfseries\color{blue},
  commentstyle=\itshape\color{purple},
  identifierstyle=\color{black},
  stringstyle=\color{red},
  extendedchars=\true,
}

%--------------------------------------
%			НАЧАЛО ДОКУМЕНТА
%--------------------------------------

\begin{document}

%--------------------------------------
%			ТИТУЛЬНЫЙ ЛИСТ
%--------------------------------------
\begin{titlepage}
\thispagestyle{empty}
\newpage


%Шапка титульного листа
%--------------------------------------
\vspace*{-60pt}
\hspace{-65pt}
\begin{minipage}{0.3\textwidth}
\hspace*{-20pt}\centering

\includegraphics[width=\textwidth]{emblem}
\end{minipage}
\begin{minipage}{0.67\textwidth}\small \textbf{
\vspace*{-0.7ex}
\hspace*{-6pt}\centerline{Министерство науки и высшего образования Российской Федерации}
\vspace*{-0.7ex}
\centerline{Федеральное государственное бюджетное образовательное учреждение }
\vspace*{-0.7ex}
\centerline{высшего образования}
\vspace*{-0.7ex}
\centerline{<<Московский государственный технический университет}
\vspace*{-0.7ex}
\centerline{имени Н.Э. Баумана}
\vspace*{-0.7ex}
\centerline{(национальный исследовательский университет)>>}
\vspace*{-0.7ex}
\centerline{(МГТУ им. Н.Э. Баумана)}}
\end{minipage}
%--------------------------------------

%Полосы
%--------------------------------------
\vspace{-25pt}
\hspace{-35pt}\rule{\textwidth}{2.3pt}

\vspace*{-20.3pt}
\hspace{-35pt}\rule{\textwidth}{0.4pt}
%--------------------------------------

\vspace{1.5ex}
\hspace{-35pt} \noindent \small ФАКУЛЬТЕТ\hspace{80pt} <<Информатика и системы управления>>

\vspace*{-16pt}
\hspace{47pt}\rule{0.83\textwidth}{0.4pt}

\vspace{0.5ex}
\hspace{-35pt} \noindent \small КАФЕДРА\hspace{50pt} <<Теоретическая информатика и компьютерные технологии>>

\vspace*{-16pt}
\hspace{30pt}\rule{0.866\textwidth}{0.4pt}
  
\vspace{11em}

\begin{center}
\Large {\bf Лабораторная работа № 8} \\ 
\large {\bf<<Изучение эффективности модифицированных методов быстрого умножения матриц>>} \\ 
\large {\bf по курсу <<Численные методы линейной алгебры>>} \\ 
\end{center}\normalsize

\vspace{8em}


\begin{flushright}
  {Студент группы ИУ9-71Б Лысенко В. В.\hspace*{15pt} \\
  \vspace{2ex}
  Преподаватель Посевин Д. П.\hspace*{15pt}}
\end{flushright}

\bigskip

\vfill
 

\begin{center}
\textsl{Москва 2024}
\end{center}
\end{titlepage}
%--------------------------------------
%		КОНЕЦ ТИТУЛЬНОГО ЛИСТА
%--------------------------------------

\renewcommand{\ttdefault}{pcr}

\setlength{\tabcolsep}{3pt}
\newpage
\setcounter{page}{2}

\section{Цель работы}\label{Sect::task}
\par
\renewcommand{\labelenumii}{\arabic{enumi}.\arabic{enumii}.}
Реализовать и исследовать эффективности модификаций методов Штрассена и Винограда-Штрассена быстрого умножения матриц.

\section{Задачи}\label{Sect::task}
\par
\renewcommand{\labelenumii}{\arabic{enumi}.\arabic{enumii}.}
\begin{enumerate} 
    \item На основе уже реализованного в лабораторной работе №7 метода Штрассена необходимо выполнить его многопоточную версию.
    \item Реализовать классическую версию метода Винограда-Штрассена.
    \item Реализовать многопоточную версию метода Винограда-Штрассена.
    \item Сравнить скорости умножения матриц всеми реализованными методами: в лабораторной работе №7 и лабораторной работе №8.
    \item Построить на одной координатной плоскости разными цветами графики зависимости времени Т (мс) умножения двух матриц размера N x N стандартным алгоритмом, алгоритмом Винограда, методом Штрассена, многопоточным методом Штрассена, методом Винограда-Штрассена, многопоточным методом Винограда-Штрассена от размера матрицы N. N изменяется от 2 до 1024. Выбор значения размера матрицы N необходимо осуществлять по формуле N=$2^n$, при этом изменяя n от 1 до 10 строго. Во время построения графиков на координатной плоскости на горизонтальной оси время, которое затрачено на выполнение алгоритма перемножения матриц строго в миллисекундах.
\end{enumerate}

\pagebreak
\section{Практическая реализация}

\par
Код программы:
\begin{lstlisting}
using LinearAlgebra
using Plots
using Base.Threads

function standard(A, B)
    n = size(A, 1)
    C = zeros(Float64, n, n)
    for i in 1:n
        for j in 1:n
            for k in 1:n
                C[i, j] += A[i, k] * B[k, j]
            end
        end
    end
    return C
end


function winograd(A, B)
    n = size(A, 1)
    C = zeros(Float64, n, n)
    m = n ÷ 2

    for i in 1:n
        for j in 1:n
            for k in 1:m
                C[i,j] += (A[i, 2*k-1] + B[2*k,j]) * (B[2*k-1,j] + A[i, 2*k]) - A[i, 2*k-1] * A[i, 2*k] - B[2*k,j] * B[2*k-1,j]

            end
            
        end
    end
    return C
end


"""
function strassen_false(A, B)
    n = size(A, 1)
    if n <= 2
        return standard(A, B)
    end
    if n % 2 != 0
        new_size = n + 1
        A_new = zeros(Float64, new_size, new_size)
        B_new = zeros(Float64, new_size, new_size)
        A_new[1:n, 1:n] = A
        B_new[1:n, 1:n] = B
        C_new = strassen_false(A_new, B_new)
        return C_new[1:n, 1:n]
    end

    k = n ÷ 2

    A11 = A[1:k, 1:k]
    A12 = A[1:k, k+1:n]
    A21 = A[k+1:n, 1:k]
    A22 = A[k+1:n, k+1:n]

    B11 = B[1:k, 1:k]
    B12 = B[1:k, k+1:n]
    B21 = B[k+1:n, 1:k]
    B22 = B[k+1:n, k+1:n]

    P1 = strassen_false(A11 + A22, B11 + B22)
    P2 = strassen_false(A21 + A22, B11)
    P3 = strassen_false(A11, B12 - B22)
    P4 = strassen_false(A22, B21 - B11)
    P5 = strassen_false(A11 + A12, B22)
    P6 = strassen_false(A21 - A11, B11 + B12)
    P7 = strassen_false(A12 - A22, B21 + B22)

    C11 = P1 + P4 - P5 + P7
    C12 = P3 + P5
    C21 = P2 + P4
    C22 = P1 + P3 - P2 + P6

    C = zeros(Float64, n, n)
    C[1:k, 1:k] = C11
    C[1:k, k+1:n] = C12
    C[k+1:n, 1:k] = C21
    C[k+1:n, k+1:n] = C22

    return C
end
"""

function strassen(A, B)
    n = size(A, 1)
    if n <= 64
        return standard(A, B)
    end
    if n % 2 != 0
        new_size = n + 1
        A_new = zeros(Float64, new_size, new_size)
        B_new = zeros(Float64, new_size, new_size)
        A_new[1:n, 1:n] = A
        B_new[1:n, 1:n] = B
        C_new = strassen(A_new, B_new)
        return C_new[1:n, 1:n]
    end

    k = n ÷ 2

    A11 = A[1:k, 1:k]
    A12 = A[1:k, k+1:n]
    A21 = A[k+1:n, 1:k]
    A22 = A[k+1:n, k+1:n]

    B11 = B[1:k, 1:k]
    B12 = B[1:k, k+1:n]
    B21 = B[k+1:n, 1:k]
    B22 = B[k+1:n, k+1:n]

    P1 = strassen(A11 + A22, B11 + B22)
    P2 = strassen(A21 + A22, B11)
    P3 = strassen(A11, B12 - B22)
    P4 = strassen(A22, B21 - B11)
    P5 = strassen(A11 + A12, B22)
    P6 = strassen(A21 - A11, B11 + B12)
    P7 = strassen(A12 - A22, B21 + B22)

    C11 = P1 + P4 - P5 + P7
    C12 = P3 + P5
    C21 = P2 + P4
    C22 = P1 + P3 - P2 + P6

    C = zeros(Float64, n, n)
    C[1:k, 1:k] = C11
    C[1:k, k+1:n] = C12
    C[k+1:n, 1:k] = C21
    C[k+1:n, k+1:n] = C22

    return C
end

function strassen_winograd(A, B)
    n = size(A, 1)
    if n <= 64
        return standard(A, B)
    end
    if n % 2 != 0
        new_size = n + 1
        A_new = zeros(Float64, new_size, new_size)
        B_new = zeros(Float64, new_size, new_size)
        A_new[1:n, 1:n] = A
        B_new[1:n, 1:n] = B
        C_new = strassen_winograd(A_new, B_new)
        return C_new[1:n, 1:n]
    end

    k = n ÷ 2

    A11 = A[1:k, 1:k]
    A12 = A[1:k, k+1:n]
    A21 = A[k+1:n, 1:k]
    A22 = A[k+1:n, k+1:n]

    B11 = B[1:k, 1:k]
    B12 = B[1:k, k+1:n]
    B21 = B[k+1:n, 1:k]
    B22 = B[k+1:n, k+1:n]

    S1 = A21 + A22
    S2 = S1 - A11
    S3 = A11 - A21
    S4 = A12 - S2
    S5 = B12 - B11
    S6 = B22 - S5
    S7 = B22 - B12
    S8 = S6 - B21

    P1 = strassen_winograd(S2, S6)
    P2 = strassen_winograd(A11, B11)
    P3 = strassen_winograd(A12, B21)
    P4 = strassen_winograd(S3, S7)
    P5 = strassen_winograd(S1, S5)
    P6 = strassen_winograd(S4, B22)
    P7 = strassen_winograd(A22, S8)

    T1 = P1 + P2
    T2 = T1 + P4

    C11 = P2 + P3
    C12 = T1 + P5 + P6
    C21 = T2 - P7
    C22 = T2 + P5

    C = zeros(Float64, n, n)
    C[1:k, 1:k] = C11
    C[1:k, k+1:n] = C12
    C[k+1:n, 1:k] = C21
    C[k+1:n, k+1:n] = C22

    return C
end


function multi_treaded_strassen(A, B)
    n = size(A, 1)
    if n <= 64
        return standard(A, B)
    end
    if n % 2 != 0
        new_size = n + 1
        A_new = zeros(Float64, new_size, new_size)
        B_new = zeros(Float64, new_size, new_size)
        A_new[1:n, 1:n] = A
        B_new[1:n, 1:n] = B
        C_new = multi_treaded_strassen(A_new, B_new)
        return C_new[1:n, 1:n]
    end

    k = n ÷ 2
    A11, A12 = A[1:k, 1:k], A[1:k, k+1:end]
    A21, A22 = A[k+1:end, 1:k], A[k+1:end, k+1:end]
    B11, B12 = B[1:k, 1:k], B[1:k, k+1:end]
    B21, B22 = B[k+1:end, 1:k], B[k+1:end, k+1:end]

    P1, P2, P3, P4, P5, P6, P7 = Tuple{Matrix{Float64}, Matrix{Float64}, Matrix{Float64}, Matrix{Float64}, Matrix{Float64}, Matrix{Float64}, Matrix{Float64}}(zeros(k, k) for _ in 1:7)

    Threads.@threads for i in 1:7
        if i == 1
            P1 .= multi_treaded_strassen(A11 + A22, B11 + B22)
        elseif i == 2
            P2 .= multi_treaded_strassen(A21 + A22, B11)
        elseif i == 3
            P3 .= multi_treaded_strassen(A11, B12 - B22)
        elseif i == 4
            P4 .= multi_treaded_strassen(A22, B21 - B11)
        elseif i == 5
            P5 .= multi_treaded_strassen(A11 + A12, B22)
        elseif i == 6
            P6 .= multi_treaded_strassen(A21 - A11, B11 + B12)
        elseif i == 7
            P7 .= multi_treaded_strassen(A12 - A22, B21 + B22)
        end
    end

    C11 = P1 + P4 - P5 + P7
    C12 = P3 + P5
    C21 = P2 + P4
    C22 = P1 - P2 + P3 + P6

    C = [C11 C12; C21 C22]

    return C[1:size(A, 1), 1:size(B, 2)]
end



function multi_treaded_strassen_winograd(A, B)
    n = size(A, 1)
    if n <= 64
        return standard(A, B)
    end
    if n % 2 != 0
        new_size = n + 1
        A_new = zeros(Float64, new_size, new_size)
        B_new = zeros(Float64, new_size, new_size)
        A_new[1:n, 1:n] = A
        B_new[1:n, 1:n] = B
        C_new = multi_treaded_strassen_winograd(A_new, B_new)
        return C_new[1:n, 1:n]
    end

    k = n ÷ 2

    A11 = A[1:k, 1:k]
    A12 = A[1:k, k+1:n]
    A21 = A[k+1:n, 1:k]
    A22 = A[k+1:n, k+1:n]

    B11 = B[1:k, 1:k]
    B12 = B[1:k, k+1:n]
    B21 = B[k+1:n, 1:k]
    B22 = B[k+1:n, k+1:n]

    S1 = A21 + A22
    S2 = S1 - A11
    S3 = A11 - A21
    S4 = A12 - S2
    S5 = B12 - B11
    S6 = B22 - S5
    S7 = B22 - B12
    S8 = S6 - B21

    P1, P2, P3, P4, P5, P6, P7 = Tuple{Matrix{Float64}, Matrix{Float64}, Matrix{Float64}, Matrix{Float64}, Matrix{Float64}, Matrix{Float64}, Matrix{Float64}}(zeros(k, k) for _ in 1:7)
    
    Threads.@threads for i in 1:7
        if i == 1
            P1 .= multi_treaded_strassen_winograd(S2, S6)
        elseif i == 2
            P2 .= multi_treaded_strassen_winograd(A11, B11)
        elseif i == 3
            P3 .= multi_treaded_strassen_winograd(A12, B21)
        elseif i == 4
            P4 .= multi_treaded_strassen_winograd(S3, S7)
        elseif i == 5
            P5 .= multi_treaded_strassen_winograd(S1, S5)
        elseif i == 6
            P6 .= multi_treaded_strassen_winograd(S4, B22)
        elseif i == 7
            P7 .= multi_treaded_strassen_winograd(A22, S8)
        end
    end

    T1 = P1 + P2
    T2 = T1 + P4

    C11 = P2 + P3
    C12 = T1 + P5 + P6
    C21 = T2 - P7
    C22 = T2 + P5

    C = zeros(Float64, n, n)
    C[1:k, 1:k] = C11
    C[1:k, k+1:n] = C12
    C[k+1:n, 1:k] = C21
    C[k+1:n, k+1:n] = C22

    return C
end

n = 100
A = randn(n, n)
B = randn(n, n)

standard_arr = standard(A, B)
winograd_arr = winograd(A, B)
strassen_arr = strassen(A, B)
strassen_winograd_arr = strassen_winograd(A, B)
multi_treaded_strassen_arr = multi_treaded_strassen(A, B)
multi_treaded_strassen_winograd_arr = multi_treaded_strassen_winograd(A, B)


error_winograd = maximum(abs.(standard_arr .- winograd_arr))
error_strassen = maximum(abs.(standard_arr .- strassen_arr))
error_strassen_winograd = maximum(abs.(standard_arr .- strassen_winograd_arr))
error_multi_treaded_strassen = maximum(abs.(standard_arr .- multi_treaded_strassen_arr))
error_multi_treaded_strassen_winograd = maximum(abs.(standard_arr .- multi_treaded_strassen_winograd_arr))

println("Погрешность алгоритма Винограда")
println(error_winograd)
println("Погрешность алгоритма Штрассена")
println(error_strassen)
println("Погрешность алгоритма Штрассена-Винограда")
println(error_strassen_winograd)
println("Погрешность алгоритма многопоточного Штрассена")
println(error_multi_treaded_strassen)
println("Погрешность алгоритма многопоточного Штрассена-Винограда")
println(error_multi_treaded_strassen_winograd)

N = [2^n for n in 1:10]
standard_times = Float64[]
winograd_times = Float64[]
strassen_times = Float64[]
strassen_winograd_times = Float64[]
multi_treaded_strassen_times = Float64[]
multi_treaded_strassen_winograd_times = Float64[]

for n in N
    A = randn(n, n)
    B = randn(n, n)
    print(n)
    println("x$n matrix counting...")

    standard_t = @elapsed standard(A, B)
    push!(standard_times, standard_t)

    winograd_t = @elapsed winograd(A, B)
    push!(winograd_times, winograd_t)

    strassen_t = @elapsed strassen(A, B)
    push!(strassen_times, strassen_t)
    
    strassen_winograd_t = @elapsed strassen_winograd(A, B)
    push!(strassen_winograd_times, strassen_winograd_t)

    multi_treaded_strassen_t = @elapsed multi_treaded_strassen(A, B)
    push!(multi_treaded_strassen_times, multi_treaded_strassen_t)
    
    multi_treaded_strassen_winograd_t = @elapsed multi_treaded_strassen_winograd(A, B)
    push!(multi_treaded_strassen_winograd_times, multi_treaded_strassen_winograd_t)
end

plot(N, standard_times, label="Стандартный алгоритм", xlabel="Размер матрицы N", ylabel="Время (с)", legend=:topleft)
plot!(N, winograd_times, label="Алгоритм Винограда")
plot!(N, strassen_times, label="Метод Штрассена")
plot!(N, strassen_winograd_times, label="Метод Штрассена-Винограда")
plot!(N, multi_treaded_strassen_times, label="Метод Штрассена (многопоточный)")
plot!(N, multi_treaded_strassen_winograd_times, label="Метод Штрассена-Винограда (многопоточный)")
\end{lstlisting}

\section{Результаты}
\par
Результат выполнения программы:
\begin{center}
\includegraphics[scale=0.8]{g1}
\end{center}
\par

\begin{center}
\includegraphics[scale=0.6]{g2}
\end{center}
\par

\section{Выводы}
В ходе выполнения лаборатороной работы были реализованы и исследованы эффективности модификаций методов Штрассена и Винограда-Штрассена быстрого умножения матриц. Резултаты выполнения показали, что в среднем сымые лучшие результаты показывает метод Штрассена-Винограда многопоточный, затем метод Штрассена-Винограда и метод Штрассена многопоточный, однако на размере матриц 1024х1024 разница не столь существенна и результаты могут меняться. Возможно на более больших размерах результаты будут более конкретные.
\end{document}
